%  LaTeX support: latex@mdpi.com 
%  For support, please attach all files needed for compiling as well as the log file, and specify your operating system, LaTeX version, and LaTeX editor.
%
%  manuscript_v3.tex: internal bibliography (MDPI template variant B, ACS format for Mathematics).
%  References are embedded below; biblio.bib is not required at compile time.
%  To refresh entries after citation changes: compile v2 with BibTeX and copy manuscript_v2.bbl.

%=================================================================
\documentclass[mathematics,article,submit,moreauthors]{Definitions/mdpi} 

%=================================================================
% MDPI internal commands - do not modify
\firstpage{1} 
\makeatletter 
\setcounter{page}{\@firstpage} 
\makeatother
\pubvolume{1}
\issuenum{1}
\articlenumber{0}
\pubyear{2026}
\copyrightyear{2026} 
%\externaleditor{Firstname Lastname} % More than 1 editor, please add `` and '' before the last editor name
\datereceived{ } 
\daterevised{ } % Comment out if no revised date
\dateaccepted{ } 
\datepublished{ } 
%\datecorrected{} % For corrected papers: "Corrected: XXX" date in the original paper.
%\dateretracted{} % For retracted papers: "Retracted: XXX" date in the original paper.
%\doinum{} % Used for some special journals, like molbank
%\pdfoutput=1 % Uncommented for upload to arXiv.org
%\CorrStatement{yes}  % For updates
%\longauthorlist{yes} % For many authors that exceed the left citation part
%\IsAssociation{yes} % For association journals

%=================================================================
% Add packages and commands here. The following packages are loaded in our class file: fontenc, inputenc, calc, indentfirst, fancyhdr, graphicx, epstopdf, lastpage, ifthen, float, amsmath, amssymb, lineno, setspace, enumitem, mathpazo, booktabs, titlesec, etoolbox, tabto, xcolor, colortbl, soul, multirow, microtype, tikz, totcount, changepage, attrib, upgreek, array, tabularx, pbox, ragged2e, tocloft, marginnote, marginfix, enotez, amsthm, natbib, hyperref, cleveref, scrextend, url, geometry, newfloat, caption, draftwatermark, seqsplit
% cleveref: load \crefname definitions after \begin{document}

\usepackage{xcolor}

%=================================================================
% Please use the following mathematics environments: Theorem, Lemma, Corollary, Proposition, Characterization, Property, Problem, Example, ExamplesandDefinitions, Hypothesis, Remark, Definition, Notation, Assumption
%% For proofs, please use the proof environment (the amsthm package is loaded by the MDPI class).

%=================================================================
% Full title of the paper (Capitalized)
%\Title{On the impact of aggregation operators and imprecise information in ranking stability: the case of prioritization problems in Cooperative Automated Decision Making}

\Title{Aggregation Semantics for Prioritization in Cooperative Automated Decision Making}

% Author Orchid ID: enter ID or remove command 
\newcommand{\orcidauthorA}{0000-0003-3267-6753} % Add \orcidA{} behind the author's name
\newcommand{\orcidauthorB}{0009-0002-0183-7905} % Add \orcidB{} behind the author's name
\newcommand{\orcidauthorC}{0000-0002-7653-1452} % Add \orcidB{} behind the author's name

% Authors, for the paper (add full first names)
\Author{Pavel Novoa-Hern\'andez$^{1,*}$\orcidA{}, Mariia Godz$^{2}$\orcidA{} and David A. Pelta$^{2}$\orcidC{}}

%\longauthorlist{yes}

% MDPI internal command: Authors, for metadata in PDF
\AuthorNames{Pavel Novoa-Hern\'andez, Mariia Godz and David A. Pelta}

%\longauthorlist{yes}

% Affiliations / Addresses (Add [1] after \address if there is only one affiliation.)
\address{%
$^{1}$ \quad Dept. Computer Engineering and Systems, Universidad de La Laguna; pnovoahe@ull.edu.es\\
$^{2}$ \quad Dept. Computer Science and Artificial Intelligence, Universidad de Granada; \{mariiagodz,dpelta\}@ugr.es}

% Contact information of the corresponding author
\corres{Correspondence:pnovoahe@ull.edu.es}

% Current address and/or shared authorship
%\firstnote{Current address: Affiliation.}  
% Current address should not be the same as any items in the Affiliation section.

%\secondnote{These authors contributed equally to this work.}
% The commands \thirdnote{} till \eighthnote{} are available for further notes.

%\simplesumm{} % Simple summary

%\conference{} % An extended version of a conference paper

% Abstract (Do not insert blank lines, i.e. \\) 
\abstract{{\color{blue}Hybrid automated decision-making systems increasingly combine predictive outputs from machine-learning models with contextual knowledge, organizational policies, or domain constraints. Within the Cooperative Automated Decision-Making System (CADEMAS) framework, the aggregation operator is not merely a mathematical component but the mechanism that defines how predictive evidence and contextual information are reconciled into a final prioritization. This paper investigates the influence of aggregation semantics by comparing the standard linear operator with two well-established alternatives: the conjunctive minimum and the weighted geometric mean. We first derive theoretical conditions characterizing rank reversals and predictive rank preservation under different aggregation regimes. These properties are then examined through Monte Carlo simulations covering policy compliance, predictive overconfidence, and uncertainty in contextual assessments. The results show that linear aggregation provides a flexible trade-off between prediction and context, but may prioritize contextually excluded alternatives when predictive evidence dominates. In contrast, the non-linear operators better preserve contextual constraints and exhibit greater robustness to predictive overconfidence and contextual uncertainty. A realistic employee attrition case study based on a CADEMAS-ML instantiation confirms that these theoretical and simulation results remain visible in practice. Overall, the findings demonstrate that the choice of aggregation operator constitutes a governance decision that directly determines the decision semantics of cooperative automated decision-making systems.}}

% Keywords
\keyword{
cooperative automated decision making;
aggregation operators;
%fuzzy aggregation;
multi-criteria decision making;
context-aware decision support;
prioritization
}

%%%%%%%%%%%%%%%%%%%%%%%%%%%%%%%%%%%%%%%%%%
\begin{document}


\section{Introduction}

%Automated Decision-Making (ADM) systems are now used in settings where recommendations have economic, organizational, and ethical consequences. 
Artificial Intelligence and machine learning techniques, in particular, have improved the ability to estimate risks and anticipate events, in many decision making scenarios. 
However, a good prediction is not always an actionable or good decision. A hospital, a company, or a public agency may need to respect resource limits, institutional policies, safety constraints, or expert-defined priorities before acting on a model output. This is why many decision-support systems are moving from purely predictive pipelines toward hybrid architectures that combine data-driven evidence with rules, contextual information, and domain knowledge \citep{Kierner2023,Rahman2025,PerezDominguez2024}.

This shift creates a methodological problem that is easy to overlook. Once predictive evidence and contextual knowledge coexist, the final recommendation depends on how both signals are integrated. The aggregation rule therefore becomes part of the decision policy. It decides whether a high predicted risk or confidence can compensate for weak contextual support, whether a contextual veto must dominate the model, or whether the system should operate somewhere between those two extremes.


{\color{blue} Recently, \citet{NovoaHernandez2026} proposed the Cooperative Automated Decision-Making System (CADEMAS) framework to make this separation explicit. Here, an automated decision-making (ADM) system is regarded as ``a (computational) process, including AI techniques and approaches, that, fed by inputs and data received or collected from the environment, can generate, given a set of predefined objectives, outputs in a wide variety of forms (content, ratings, recommendations, decisions, predictions, etc.)'' \cite{ELI_ADM_Principles_2022}. CADEMAS organizes heterogeneous ADM units within a layered architecture: predictive or automated components produce partial evidence, an integration layer combines their outputs, and a contextual modulator evaluates the result under organizational or institutional criteria.}

CADEMAS examples are those presented in \citep{NovoaHernandez2026, Godz2026}. In both cases, the system uses machine learning based ADM with fuzzy rule based context representation. In the former, the aim is to prioritize a set of small and medium enterprises to rescue, while in the later the aim is to prioritize retention strategies for avoiding employees in attrition risk. In this case, 
fuzzy contextual rules encode retention priorities such as strategic fit, development potential, or organizational vulnerability.

%In the machine-learning-based CADEMAS (CADEMAS-ML) employee-attrition instantiation of \citet{Godz2026}, for example, machine-learning models estimate the risk that an employee will leave, while fuzzy contextual rules encode retention priorities such as strategic fit, development potential, or organizational vulnerability.

These cases combine predictive score and context through a linear convex rule. The appeal is clear: the parameter is interpretable, and the final score is easy to explain. Yet linearity also imposes a strong assumption. It treats prediction or scoring and context as fully compensatory dimensions. For instance, a sufficiently high risk score may offset poor contextual compatibility, and strong contextual support may offset weak predictive evidence. This is not necessarily wrong, but it is not neutral. In domains where contextual information represents a regulation, a safety requirement, or a fairness constraint, unrestricted compensation may lead to recommendations that are numerically plausible but institutionally unacceptable.

The Multi-Criteria Decision-Making (MCDM) literature, including fuzzy MCDM, offers many alternatives to such basic linear aggregation mechanism. Ordered weighted \citep{Yager1988}, geometric \citep{Tan2009}, t-norm-based \citep{Grabisch1996}, Frank \citep{Riaz2022}, Einstein \citep{Iampan2021}, Aczel--Alsina \citep{Senapati2022}, Dombi \citep{Akram2020}, and Diophantine fuzzy \citep{Farid2022,Jeevitha2023} operators have been developed to represent different forms of interaction among criteria. 
{\color{blue}
This literature is mathematically rich and methodologically diverse \citep{Mardani2018,Sun2018,Kahraman2015,Mardani2015}. It characterizes aggregation operators in terms of formal properties, behavioral patterns, extension mechanisms, and parameterized forms that govern how criteria interactions are modeled \citep{Sun2018}. It also shows that the choice of aggregation operator significantly affect criteria weights and the final ranking of alternatives in applied decision problems \citep{Petrovic2023}.
}

Despite this extensive body of work, aggregation operators have been studied primarily as mathematical constructs or as components of individual decision-making methods. Comparatively less attention has been paid to their role within cooperative automated decision-making systems, where aggregation is responsible for reconciling the outputs of heterogeneous ADM units with contextual or organizational knowledge. In such settings, the aggregation operator is not merely a numerical device; it determines how different sources of evidence interact and therefore embodies an implicit decision policy.

This perspective is particularly relevant for the CADEMAS framework. By explicitly separating evidence generation from contextual assessment, CADEMAS allows different aggregation strategies to be applied without modifying the underlying ADM units or the contextual knowledge itself. Consequently, the aggregation mechanism becomes a key design choice: different operators may produce substantially different recommendations or decisions from exactly the same inputs. A highly compensatory strategy may allow strong evidence from one source to outweigh weak contextual support, whereas a more restrictive strategy may preserve contextual constraints even when the automated evidence is highly favorable. The choice of aggregation therefore influences not only numerical scores but also the transparency, robustness, and governance characteristics of the overall decision process.

Surprisingly, this question has received relatively little systematic attention. Although the fuzzy MCDM literature offers a rich collection of aggregation operators and analyzes many of their mathematical properties, most studies focus either on proposing new operators or on demonstrating their usefulness in specific application domains. Comparatively few works investigate how alternative aggregation semantics affect the behavior of an entire cooperative decision-making framework. Questions such as the robustness of prioritization under contextual uncertainty, the susceptibility of recommendations to localized prediction errors, or the consistency of resource-allocation policies across different aggregation strategies remain largely unexplored.

This paper addresses this gap in the context of CADEMAS, using as example a prioritization task. Rather than proposing yet another aggregation operator, we investigate how representative aggregation strategies with different compensation semantics influence the behavior of the same cooperative decision-making architecture. Our objective is to understand how the choice of aggregation affects prioritization outcomes, policy compliance, robustness to predictive perturbations, and ranking stability, while preserving the modular philosophy of CADEMAS, in which evidence generation and decision integration remain conceptually independent.

To this end, we combine controlled simulation experiments with the prior CADEMAS example on  employee-attrition introduced by \citet{Godz2026}. The simulation study isolates the effects of aggregation under different operating conditions, whereas the real case study illustrates how these behaviors translate into a practical prioritization scenario. Together, these complementary analyses provide both controlled evidence and application-oriented insight into the implications of aggregation design choices.

The main contributions of this work are as follows:

\begin{itemize}
\item We reinterpret aggregation within CADEMAS as a policy-design mechanism governing the interaction between the outputs of heterogeneous ADM units and contextual knowledge.

\item We compare representative aggregation regimes exhibiting different compensation semantics within a common cooperative decision-making framework.

\item We provide a systematic experimental evaluation of how aggregation semantics affect prioritization behavior, policy compliance, robustness to localized predictive perturbations, and ranking stability.

\item We demonstrate these effects through controlled simulation and a CADEMAS-ML employee-attrition case study, highlighting practical implications of aggregation design choices.

\end{itemize}

The remainder of the paper is organized as follows. Section~\ref{sect:cademas} introduces the CADEMAS framework and its instantiations. Section~\ref{sect:aggregation_operators} formalizes the aggregation models and their structural properties. Section~\ref{sect:simulation} presents the simulation study and its findings. Section~\ref{sect:real_example} applies the operators to a published employee-attrition instantiation. Section~\ref{sect:conclusions} summarizes the main conclusions, limitations, and directions for future work.

\section{The CADEMAS Framework and Its Instantiations}
\label{sect:cademas} 

%\textbf{EN ALGUN LADO TENDRIAMOS QUE MENCIONAR ALGO MÁS DEL CONTEXTO Y CITAR NUESTRO PAPER}

\citet{NovoaHernandez2026} introduced CADEMAS as a general framework for designing context-aware decision-support systems in which multiple heterogeneous Automated Decision-Making (ADM) units cooperate to produce a unified recommendation, score or label. The framework is intentionally independent of the internal nature of these decision units: an ADM may correspond to a predictive model, a rule-based system, an optimization algorithm, an expert system, or any other computational mechanism capable of producing decision evidence. 
Likewise, contextual knowledge may be represented through expert rules, organizational policies, regulatory constraints, or alternative knowledge-representation formalisms.

Rather than prescribing how individual decisions are obtained, CADEMAS defines an architectural process through which heterogeneous decision signals are coordinated and subsequently governed by \textit{contextual} knowledge before a final recommendation is produced.

Formally, a CADEMAS instance is represented by the tuple
\begin{equation}
\mathcal{S}=
\langle I, U, O, K, T \rangle,
\end{equation}
where $I$ denotes the Input Manager, $U$ the set of ADM units, $O$ the Decision Integration Layer, $K$ the Contextual Modulator, and $T$ the universe of admissible data representations. %The framework intentionally abstracts away from the internal implementation of these components. %Consequently, an ADM unit may correspond to a machine-learning model, a rule-based system, an optimization module, or any other computational mechanism capable of producing decision evidence. Likewise, contextual knowledge may be represented through expert rules, organizational policies, regulatory constraints, or alternative knowledge-representation formalisms.

The resulting architecture follows a layered decision process. First, the ADM units independently generate decision evidence from the available input information. Note that every ADM may use different parts of that information. 
The Contextual Modulator, using again all or a part of the input information, produces a ``degree of satisfaction'' with the domain knowledge, organizational priorities, regulatory requirements, or other contextual considerations. 
Finally, the Decision Integration Layer combines the predictive evidence with the contextual's satisfaction degree into a unified decision score. 

This separation between evidence generation and contextual governance is one of the defining characteristics of CADEMAS and enables the same architecture to be instantiated across diverse application domains without modifying its overall decision process.

One important specialization of this framework is the so called CADEMAS-ML, in which the ADM units correspond to machine-learning predictive models. In this case, the predictive subsystem resembles an ensemble architecture in which multiple predictive models cooperate to produce a unified predictive assessment. However, unlike conventional ensemble methods, CADEMAS-ML explicitly separates predictive cooperation from contextual governance. Predictive evidence is first consolidated into a cooperative predictive score and only afterwards, reconciled with \textit{fuzzy} contextual knowledge that encodes organizational or institutional priorities. {\color{blue} This contextual knowledge is represented through a fuzzy inference system inspired by  \citet{PerezCanedo2022} and \citet{PerezCanedo2024}, enabling organizational and institutional priorities to be incorporated in an interpretable manner.}

For better illustrating the main ideas of this contribution, in what follows we will focus exclusively on CADEMAS-ML. Independently of the application domain, we abstract the predictive subsystem by a normalized cooperative \emph{predictive score}
\begin{equation}
R_i \in [0,1],
\end{equation}
which summarizes the collective evidence generated by the cooperating ADM units for decision case $i$. The precise semantics of $R_i$ depend on the application: it may represent, for example, a cooperative risk score, a confidence score, or another normalized predictive assessment.

Independently, the contextual subsystem computes a normalized contextual satisfaction degree
\begin{equation}
Q_i \in [0,1],
\end{equation}
also referred to as the \emph{context score}, which quantifies the extent to which decision case $i$ satisfies the rules of the context.

The predictive and contextual assessments are subsequently combined through an aggregation operator $A$
\begin{equation}
P_i = A(R_i,Q_i),
\end{equation}
to obtain a final \textit{prioritization} score. As we are dealing with a prioritization example, $P_i$ will be used to rank candidate cases and obtain the corresponding Top-$K$ candidates.

Existing CADEMAS-ML applications employ a \textit{linear} convex aggregation operator to combine predictive scores and contextual information. However, the CADEMAS framework itself does not prescribe any particular aggregation semantics \citep{NovoaHernandez2026}. The integration operator can therefore be regarded as an independent design component whose behavior determines how predictive scores and contextual knowledge interact. The remainder of this paper investigates alternative aggregation operators and analyzes how their semantics reshape the resulting prioritization process while preserving the overall CADEMAS-ML architecture.

{\color{blue}
\section{Aggregation Operators for Integrating Predictive Score and Context}
\label{sect:aggregation_operators}

This section formalizes the integration stage of CADEMAS-ML, where the cooperative predictive score and the contextual assessment are combined into a single prioritization score. Rather than restricting the integration mechanism to the linear convex combination employed in previous CADEMAS-ML applications, we consider a broader family of aggregation operators inspired by fuzzy logic and Multi-Criteria Decision Making (MCDM). This general formulation enables different integration semantics to be analyzed within a common mathematical framework.

Let
\[
A:[0,1]^2 \rightarrow [0,1]
\]
denote an aggregation operator mapping the cooperative predictive score $R_i$ and the contextual assessment $Q_i$ into the final priority score
\[
P_i = A(R_i,Q_i).
\]

Throughout this work, aggregation operators are assumed to satisfy the following properties:

\begin{itemize}
\item \textit{Monotonicity:} if $R_i\le R_j$ and $Q_i\le Q_j$, then
\[
A(R_i,Q_i)\le A(R_j,Q_j).
\]

\item \textit{Boundary conditions:}
\[
A(0,0)=0,\qquad A(1,1)=1.
\]

\item \textit{Normalization:}
\[
A:[0,1]^2\rightarrow[0,1].
\]
\end{itemize}

Within this general framework, the operator $A$ may implement different integration semantics. In this paper we distinguish four representative families:

\begin{itemize}
\item $A_L$: linear (fully compensatory) aggregation;
\item $A_C$: conjunctive (non-compensatory) aggregation;
\item $A_D$: disjunctive (permissive) aggregation;
\item $A_G$: non-linear compensatory aggregation based on the weighted geometric mean.
\end{itemize}

The linear operator corresponds to the integration mechanism adopted in previous CADEMAS-ML implementations and serves as the reference model throughout this work. The remaining operators represent alternative integration policies with progressively different compensation behavior between predictive evidence and contextual information.
}

\subsection{Linear Aggregation: The Current CADEMAS-ML Integration Scheme}
\label{subsect:linear}

The linear aggregation operator,
\begin{equation}
A_L(R_i,Q_i)=\lambda R_i+(1-\lambda)Q_i,\qquad \lambda\in[0,1],
\end{equation}
corresponds to the integration mechanism adopted in current CADEMAS-ML applications. The parameter $\lambda$ determines the relative importance assigned to the cooperative predictive score ($R_i$) and the contextual assessment ($Q_i$), yielding the final priority score
\[
P_i=A_L(R_i,Q_i).
\]

The operator $A_L$ is fully compensatory. Consider two candidates $i$ and $j$ such that $R_i>R_j$ but $Q_i<Q_j$. Under linear aggregation, their relative ordering depends solely on the value of $\lambda$: any decrease in contextual assessment can always be offset by a sufficiently large increase in predictive score, and vice versa. Consequently, an alternative with weak contextual support may still be prioritized if its predictive performance is sufficiently high.

This compensatory behavior follows directly from the constant partial derivatives
\[
\frac{\partial A_L}{\partial R_i}=\lambda,\qquad
\frac{\partial A_L}{\partial Q_i}=1-\lambda,
\]
which imply a constant marginal rate of substitution between prediction and context. Equivalently, the iso-priority contours in the $(R_i,Q_i)$ plane are straight lines, reflecting the same trade-off throughout the entire score domain. As a result, contextual deficiencies remain fully compensable regardless of their severity.

Although such behavior is appropriate when prediction and context represent equally substitutable sources of information, it may be undesirable in decision scenarios where contextual information encodes safety requirements, fairness constraints, regulatory compliance, or institutional vetoes. In these settings, unrestricted compensation may assign high priority to alternatives that should instead be penalized or excluded from consideration.

{\color{blue}

\subsection{Alternative Aggregation Operators}
\label{subsect:alternative_operators}

Beyond the linear operator $A_L$, we consider aggregation operators representing different degrees of compensation between predictive evidence and contextual assessment. These operators are well established in fuzzy aggregation and MCDM, and each induces a distinct integration policy within CADEMAS-ML.

\begin{enumerate}

\item \emph{Conjunctive (non-compensatory) aggregation.}
Classical fuzzy conjunction is represented by operators such as the minimum and the algebraic product \citep{Zimmermann1980}:
\begin{equation}
A_C^{\min}(R_i,Q_i)=\min(R_i,Q_i),
\end{equation}
\begin{equation}
A_C^{\prod}(R_i,Q_i)=R_iQ_i.
\end{equation}

Both operators prevent full compensation between prediction and context. Under the minimum operator, the final priority cannot exceed the weaker component, so a poor contextual assessment directly limits the overall score. In particular, if $Q_i=0$, then $P_i=0$ regardless of the predictive score, naturally implementing an institutional veto. This interpretation is consistent with the role of non-compensatory aggregation in composite indicators and policy evaluation \citep{Natoli2024,Blancas2024}. The algebraic product shares the same zero-absorption property while providing a smoother penalization of intermediate values than the minimum.

For completeness, we also note the disjunctive operator
\begin{equation}
A_D(R_i,Q_i)=\max(R_i,Q_i),
\end{equation}
which represents the opposite integration philosophy by allowing either component to dominate the final score \citep{Zimmermann1980}. Although appropriate in exploratory or opportunity-driven decision problems, this operator is incompatible with the governance objectives of CADEMAS-ML: a case with $Q_i=0$ would still obtain $P_i=R_i$, allowing predictive evidence alone to override contextual exclusion. Consequently, the disjunctive regime is not considered further in this work.

\item \emph{Non-linear compensatory aggregation.}
An intermediate degree of compensation is provided by the weighted geometric mean,
\begin{equation}
A_G(R_i,Q_i)=R_i^{\lambda}Q_i^{1-\lambda},
\qquad \lambda\in[0,1].
\end{equation}

Unlike the linear operator, the weighted geometric mean penalizes imbalance between $R_i$ and $Q_i$, while still permitting partial compensation. Consequently, it occupies an intermediate position between the arithmetic mean and conjunctive operators: contextual deficiencies have a stronger impact on the final priority than under linear aggregation, yet they do not completely suppress it unless $Q_i=0$ \citep{Natoli2024,Blancas2024}.

\end{enumerate}
}

{\color{blue}

\subsection{Rank Formulation and Structural Properties}
\label{subsect:rank_properties}

Given a set of alternatives $X=\{1,\ldots,n\}$, each alternative $i\in X$ is characterized by a cooperative predictive score $R_i\in[0,1]$ and a contextual assessment $Q_i\in[0,1]$. For a given aggregation operator $A$, the final priority score is $P_i=A(R_i,Q_i)$.

The cooperative predictive scores induce the predictive ranking
\begin{equation}
\pi_R=\operatorname{argsort}_{i\in X}R_i,
\end{equation}
which corresponds to the ordering that would be obtained by a purely predictive CADEMAS-ML configuration, i.e., without contextual modulation.

For the same aggregation operator, let $\mathbf{P}(A)=(P_1,\ldots,P_n)$ denote the aggregated score vector and
\begin{equation}
\pi(A)=\operatorname{argsort}_{i\in X}P_i,
\end{equation}
the resulting priority ranking, such that $P_{\pi(A)(1)}\ge P_{\pi(A)(2)}\ge\cdots\ge P_{\pi(A)(n)}$.

The mapping $A \longrightarrow \mathbf{P}(A) \longrightarrow \pi(A)$
therefore represents the operational output of the CADEMAS-ML integration layer, determining which alternatives belong to the Top-$K$ intervention tier and in which order.

The following analysis characterizes how the choice of aggregation operator modifies the final ranking $\pi(A)$ relative to the predictive ranking $\pi_R$. We focus on two structural properties that determine the behavior of the integration mechanism independently of any particular application domain.

The first property concerns \emph{aggregation-induced rank reversal}, namely the situation where the relative ordering of two alternatives changes after contextual information is incorporated through the aggregation operator.

Let $i,j\in X$ be two alternatives such that
\[
R_i>R_j,
\]
meaning that the predictive layer ranks alternative $i$ ahead of alternative $j$. The following proposition characterizes the conditions under which contextual information can reverse this ordering under the different aggregation regimes considered in this work.

}

\begin{Proposition}[Policy-Driven Veto and Forced Rank Reversal]
\label{prop:rank_reversal}
The following results characterize the conditions under which the aggregation mechanism induces rank reversals under linear and non-linear regimes:

\begin{enumerate}
\item Under linear convex aggregation $A_L$ with trade-off parameter $\lambda \in (0,1)$, a rank reversal occurs ($P_j > P_i$) if and only if the contextual advantage of alternative $j$, defined as $\Delta Q = Q_j - Q_i$, satisfies:
\begin{equation}
\Delta Q > \frac{\lambda}{1-\lambda}(R_i - R_j).
\end{equation}

\item Under a strict conjunctive aggregator $A_C$ (which includes both the minimum $A_C^{\min}$ and the algebraic product $A_C^{\prod}$), if the contextual score of the disadvantaged alternative is zero ($Q_i = 0$) and the other alternative has strictly positive predictive score and context score ($R_j > 0, Q_j > 0$), then a rank reversal is guaranteed: $P_j > P_i$, independently of the magnitude of $R_i$. This holds because $A_C(R_i, 0) = 0$ and $A_C(R_j, Q_j) > 0$.

\item Under the non-linear compensatory regime modeled by the weighted geometric mean $A_G$, assuming strictly positive inputs, a rank reversal occurs ($P_j > P_i$) if and only if the logarithmic ratio of contextual advantage satisfies:
\begin{equation}
\ln\left(\frac{Q_j}{Q_i}\right) > \frac{\lambda}{1-\lambda} \ln\left(\frac{R_i}{R_j}\right).
\end{equation}
\end{enumerate}
\end{Proposition}

\begin{proof}
For part (1), we establish the equivalence by starting from the rank reversal condition:
$$P_j > P_i \iff \lambda R_j + (1-\lambda) Q_j > \lambda R_i + (1-\lambda) Q_i.$$
Rearranging the terms (a strictly reversible algebraic operation) yields:
$$\iff (1-\lambda)(Q_j - Q_i) > \lambda (R_i - R_j).$$
Given $\Delta Q = Q_j - Q_i$, and since $\lambda \in (0,1)$ ensures $(1-\lambda) > 0$, we can divide both sides by $(1-\lambda)$ without altering the inequality's direction or the logical equivalence:
$$\iff \Delta Q > \frac{\lambda}{1-\lambda}(R_i - R_j).$$
This establishes both the necessary and sufficient conditions for the rank reversal under linear aggregation.

For part (2), we note that any strict conjunctive operator $A_C$ satisfies $A_C(x,0)=0$ and $A_C(x,y)>0$ whenever $x>0$ and $y>0$. Since $Q_i=0$, we have $P_i = A_C(R_i,0)=0$. Given $R_j>0$ and $Q_j>0$, it follows that $P_j = A_C(R_j,Q_j) > 0$, hence $P_j > P_i$ regardless of the value of $R_i$.

Finally, for part (3), assuming $R_i, R_j, Q_i, Q_j > 0$, the rank reversal condition under the weighted geometric mean is:
$$P_j > P_i \iff R_j^\lambda Q_j^{1-\lambda} > R_i^\lambda Q_i^{1-\lambda}.$$
Taking the natural logarithm of both sides (a strictly monotonically increasing transformation that preserves the inequality) yields:
$$\iff \lambda \ln(R_j) + (1-\lambda) \ln(Q_j) > \lambda \ln(R_i) + (1-\lambda) \ln(Q_i).$$
Rearranging the terms gives:
$$\iff (1-\lambda)(\ln(Q_j) - \ln(Q_i)) > \lambda (\ln(R_i) - \ln(R_j)).$$
Using logarithmic properties and dividing by $(1-\lambda) > 0$, we obtain the desired equivalence:
\[
\iff \ln\left(\frac{Q_j}{Q_i}\right) > \frac{\lambda}{1-\lambda} \ln\left(\frac{R_i}{R_j}\right).
\]
This completes the proof.
\end{proof}

The first part of Proposition~\ref{prop:rank_reversal} states the compensation boundary of linear CADEMAS integration: a low contextual score $Q_i$ can always be offset by a sufficiently large predictive score differential. Contextual signals then act as soft adjustments rather than binding constraints. The second part encodes the opposite logic. Under conjunctive operators, $Q_i = 0$ is an institutional veto that no amount of predictive evidence can overcome. The case is excluded from prioritization regardless of $R_i$, which protects policy compliance but may displace a high-predictive-score alternative in the ranking.

The third part describes the geometric mean as an intermediate regime. The reversal threshold depends on score ratios rather than score differences, so compensation is asymmetric and scale-dependent. In CADEMAS terms, small contextual advantages can outweigh large predictive-score differentials when both scores are low, whereas high scores require proportionally larger contextual gains to reverse the ranking. This offers a softer governance option between unrestricted linear trade-offs and hard veto semantics.

Together, these three regimes form a spectrum of integration policies that CADEMAS designers can encode through operator choice: fully compensatory ($A_L$), strictly veto-based ($A_C^{\min}$ or $A_C^{\prod}$), or non-linear compensatory with a zero boundary ($A_G$).\\

The second structural property addresses the complementary situation where contextual conditions are homogeneous across alternatives.

\begin{Proposition}[Contextual Homogeneity and Predictive Rank Preservation]
\label{prop:cont_homog_rank_preserv}
Let $Y \subseteq X$ be a subset of alternatives sharing a homogeneous contextual evaluation, such that $Q_i = q$ for all $i \in Y$, with $q > 0$. For any aggregation operator $A(R_i, Q_i)$ that is strictly monotonic with respect to its first argument when $Q_i > 0$, the strict ordinal ranking induced by the predictive model is fully preserved within $Y$. That is:
\begin{equation}
P_i > P_j \iff R_i > R_j \quad \forall i,j \in Y.
\end{equation}
\end{Proposition} 

\begin{proof}
Since $Q_i = Q_j = q$ for any $i,j \in Y$, the aggregated scores are $P_i = A(R_i, q)$ and $P_j = A(R_j, q)$.
The direct implication ($R_i > R_j \implies P_i > P_j$) follows directly from the strict monotonicity of $A$ with respect to its first argument (with the second fixed at $q > 0$). For the reverse implication ($P_i > P_j \implies R_i > R_j$), we proceed by contraposition: if we assume $R_i \leq R_j$, the strict monotonicity condition would require $A(R_i, q) \leq A(R_j, q)$, meaning $P_i \leq P_j$. This contradicts our initial premise. Therefore, the double implication holds: 
\[
P_i > P_j \iff R_i > R_j
\]
ensuring that the strict ordinal ranking induced by the predictive model is fully preserved.
\end{proof} 

Proposition~\ref{prop:cont_homog_rank_preserv} guarantees that fuzzy contextual integration does not arbitrarily disrupt the predictive ranking when there is no contextual reason to differentiate between cases. If all alternatives in a group receive the same context score ($Q_i = q$), the CADEMAS ranking within that group follows the predictive evidence order exactly. Contextual modulation then affects the absolute priority level but not the relative order among equally treated cases.

Together, these propositions establish the theoretical bounds for the empirical analysis presented in Section~\ref{sect:simulation}. Proposition~\ref{prop:rank_reversal} defines when context can override prediction under linear, non-linear compensatory, and veto-based regimes. Proposition~\ref{prop:cont_homog_rank_preserv} ensures that such overrides require genuine contextual differentiation rather than arising from the aggregation mechanism alone. 

\section{Experimental analysis}\label{sect:simulation}

The previous section established the theoretical properties of the aggregation operators. We now examine how these properties manifest in representative CADEMAS-ML prioritization scenarios. To this end, we conduct Monte Carlo simulations on synthetic populations designed to isolate three practical aspects of the integration process:

\begin{enumerate}

\item \emph{Policy compliance.} How does increasing the influence of the predictive score affect the ability of each operator to prevent contextually excluded alternatives ($Q_i=0$) from entering the Top-$K$ prioritized tier?

\item \emph{Robustness to predictive overconfidence.} Can the integration process prevent artificially inflated predictive scores from promoting contextually excluded alternatives?

\item \emph{Rank stability under imprecision in contextual evaluation.} How sensitive is the final prioritization to perturbations in the context score?

\end{enumerate}

We run a set of experiments considering a universe of $n=1000$ candidates/alternatives from which only the Top-$K$ tier ($K=100$) are supposed to receive intervention. Depending on the experiment, we evaluate the linear ($A_L$), minimum ($A_C^{\min}$), and weighted geometric ($A_G$) aggregation operators. The trade-off parameter $\lambda$ is either swept or fixed according to the objective of each experiment. Performance is assessed using three complementary metrics. Let $\Pi_K$ denote the selected Top-$K$ tier and $X_{\mathrm{veto}}=\{i:Q_i=0\}$ the set of contextually excluded candidates. The policy violation rate is defined as
\[
V_K=\frac{|\Pi_K\cap X_{\mathrm{veto}}|}{K},
\]
while the average predictive score is computed as
\[
\bar{R}_K=\frac{1}{|\Pi_K\setminus X_{\mathrm{veto}}|}
\sum_{i\in \Pi_K\setminus X_{\mathrm{veto}}}R_i,
\]
thus measuring the predictive quality of the selected compliant alternatives. 

Each experiment consists of $N=30$ independent Monte Carlo trials. Reported curves correspond to trial averages, while shaded regions indicate non-parametric 95\% confidence intervals obtained from the empirical percentiles. Throughout the analysis, we describe a behavior as \emph{structurally robust} when it is observed in every trial, and as \emph{statistically supported} when the corresponding confidence intervals remain consistently separated.

\begin{table}[t] 
\caption{Simulation protocol by aspect explored.}
\label{tab:sim_protocol}
\renewcommand{\arraystretch}{1.3}\footnotesize
\begin{adjustwidth}{-\extralength}{0cm}
\begin{tabularx}{\fulllength}{>{\raggedright\arraybackslash}X|C|C|C}
\toprule
\textbf{Parameter} & \textbf{Policy compliance} & \textbf{Algorithmic firewall} & \textbf{Rank stability} \\
\midrule
Universe size ($n$) & 1000 & 1000 & 1000 \\
\midrule
Top-$K$ tier size ($K$) & 100 & 100 & 100 \\
\midrule
Monte Carlo trials ($N$) & 30 & 30 & 30 \\
\midrule
Population design &
$X_{\text{std}}$ (80\%): $R_i, Q_i \sim \mathcal{U}(0,1)$;\ $X_{\text{veto}}$ (20\%): $Q_i = 0$ &
$X_{\text{std}}$ (80\%): $R_i, Q_i \sim \mathcal{U}(0,1)$;\ $X_{\text{veto}}$ (20\%): $Q_i = 0$ &
All cases: $R_i, Q_i \sim \mathcal{U}(0,1)$ \\
\midrule
Veto-group risk ($R_i$) &
$R_i \sim \mathcal{B}(8,2)$, fixed &
$R_i \sim \mathcal{N}(\mu_{\text{shift}}, 0.05^2)$, $\mu_{\text{shift}} \in [0.5, 1.0]$ &
-- \\
\midrule
Context perturbation &
-- &
-- &
$Q_i' = \mathrm{clip}(Q_i + \epsilon, 0, 1)$, $\epsilon \sim \mathcal{N}(0, \sigma^2)$, $\sigma \in [0, 0.5]$ \\
\midrule
Aggregation operators &
$A_L$, $A_G$ &
$A_L$, $A_C^{\min}$, $A_G$ &
$A_L$, $A_C^{\min}$ \\
\midrule
Trade-off parameter ($\lambda$) &
$[0.05, 0.95]$ (19 values) &
$0.85$ fixed &
$0.35$ fixed \\
\midrule
Evaluation metrics &
$V_K$, $\bar{R}_K$ &
$V_K$ &
Kendall's rank correlation - $\bar{\tau}(\sigma)$ \\
\bottomrule
\end{tabularx}
\end{adjustwidth} 
\end{table}

Table~\ref{tab:sim_protocol} summarizes the complete experimental protocol. Rather than generating complete data for the candidates, each synthetic decision population is modeled directly by normalized tuples $(R_i,Q_i)$, where $R_i$ is the predictive score and $Q_i$ the context score. Unless explicitly stated otherwise, both quantities are sampled independently from $\mathcal{U}(0,1)$, whereas contextual perturbations are drawn from $\mathcal{N}(0,\sigma^2)$. {\color{blue} In addition to the uniform setting, selected experiments consider predictive scores sampled from a Beta distribution, $R_i\sim\mathcal{B}(\alpha,\beta)$, to emulate predictive systems with systematic scoring biases. For example, $\mathcal{B}(8,2)$ (i.e., $\alpha=8$, $\beta=2$) concentrates most probability mass near one, representing scenarios in which the predictive subsystem consistently assigns high-risk (or high-priority) scores to most candidates.}

\subsection{Policy Compliance}
\label{subsect:opportunity_cost}

This experiment examines how increasing the influence of the predictive score affects policy compliance. Specifically, we analyze the trade-off between enforcing contextual exclusions ($V_K = 0$) and increasing the average predictive score of the selected Top-$K$ alternatives as the trade-off parameter $\lambda$ increases.

Figure~\ref{fig:opportunity_cost} compares the behavior of the linear ($A_L$) and weighted geometric ($A_G$) operators. Panel~(a) reports the policy violation rate ($\bar{V}_K$), whereas panel~(b) shows the average predictive score ($\bar{R}_K$) of the selected compliant alternatives.

\begin{figure}[t]
\centering
\includegraphics[width=\textwidth]{figures/opportunity_cost.pdf}
\caption{Policy compliance analysis ($N=30$ Monte Carlo trials). Solid curves represent Monte Carlo means and shaded regions denote 95\% confidence intervals. (\textbf{a}) Policy violation rate ($\bar{V}_K$) as a function of $\lambda$. (\textbf{b}) Average predictive score ($\bar{R}_K$) of the selected non-vetoed Top-$K$ alternatives.} 
\label{fig:opportunity_cost}
\end{figure}

The linear operator exhibits a clear compliance limit. As $\lambda$ increases, the average predictive score of the selected alternatives improves steadily while policy compliance is preserved. Specifically, $A_L$ maintains $\bar{V}_K=0$ up to approximately $\lambda=0.75$--$0.80$, during which $\bar{R}_K$ increases from about $0.53$ to $0.92$. Beyond this point, however, contextually excluded alternatives begin to enter the prioritized tier. At $\lambda=0.95$, approximately 25\% of the Top-$K$ alternatives belong to the veto group ($\bar{V}_K\approx0.25$, 95\% CI $[0.14,0.34]$), despite achieving an average predictive score close to $0.95$.

This behavior is consistent with Proposition~\ref{prop:rank_reversal}-(1). Under linear aggregation, sufficiently large predictive scores compensate for a zero context score, allowing contextually excluded alternatives to enter the prioritized set once the predictive score dominates the aggregation.

The weighted geometric operator behaves differently. Because $Q_i=0$ always implies $P_i=0$, policy compliance is preserved throughout the entire range of $\lambda$, with $\bar{V}_K=0$ in all Monte Carlo trials. At the same time, the average predictive score continues to increase as $\lambda$ grows, reaching values comparable to those obtained with the linear operator without admitting contextually excluded alternatives.

Overall, these results show that the limitation of linear aggregation is not its ability to exploit the predictive score at moderate values of $\lambda$, but its inability to preserve both high predictive scores and full policy compliance when prediction becomes the dominant component of the prioritization. In contrast, the weighted geometric operator extends the range of admissible $\lambda$ values over which both objectives can be achieved simultaneously.

\subsection{Robustness to Predictive Overconfidence}
\label{subsect:firewall}

This experiment evaluates how localized predictive overconfidence affects policy compliance. Specifically, predictive scores within the veto group are progressively increased while their context scores remain fixed at $Q_i=0$, allowing us to assess whether artificially inflated predictive scores can promote contextually excluded alternatives into the Top-$K$ tier. 

%{\color{red} PAVEL: hay que explicar el parámetro $\mu_{shift}$ aqui}

{\color{blue} As noted in Table~\ref{tab:sim_protocol}, predictive overconfidence is modeled through the parameter $\mu_{\text{shift}}$, which controls the mean predictive score assigned to the veto subgroup. Specifically, predictive scores for vetoed alternatives are generated from $R_i \sim \mathcal{N}(\mu_{\text{shift}},0.05^2)$,
truncated to the interval $[0,1]$. Increasing $\mu_{\text{shift}}$ therefore simulates progressively more adverse scenarios in which the predictive subsystem systematically assigns high scores to alternatives that should ultimately be rejected by the contextual layer.}

The trade-off parameter is fixed at $\lambda=0.85$, where the predictive score has sufficient influence for compensatory effects to emerge. Lower values of $\lambda$ would cap the priority score of vetoed alternatives below the empirical Top-$K$ threshold, masking the differences between aggregation operators.

Figure~\ref{fig:predictive_overconfidence} reports the policy violation rate ($\bar{V}_K$) as the mean predictive score of the veto group increases.
The linear operator remains fully compliant while the predictive scores of the veto group are moderate. However, once their average predictive score exceeds approximately $0.80$, policy violations increase steadily. At $\mu_{\text{shift}}=1.0$, nearly 30\% of the prioritized Top-$K$ alternatives belong to the veto group ($\bar{V}_K\approx0.30$, 95\% CI $[0.10,0.48]$). This behavior indicates that sufficiently large predictive scores compensate for the absence of contextual support, allowing contextually excluded alternatives to enter the prioritized set.

\begin{figure}[t]
\centering
\includegraphics[width=\textwidth]{figures/predictive_overconfidence.pdf}
\caption{Robustness to predictive overconfidence ($\lambda=0.85$, $N=30$ Monte Carlo trials). Solid curves represent Monte Carlo means and shaded regions denote 95\% confidence intervals. The horizontal axis controls the mean predictive score assigned to the veto group. The vertical axis reports the policy violation rate ($\bar{V}_K$).}
\label{fig:predictive_overconfidence}
\end{figure}

By contrast, both the minimum and weighted geometric operators maintain $\bar{V}_K=0$ throughout the entire range of predictive shifts. Since $Q_i=0$ always implies $P_i=0$, increasing the predictive score alone cannot promote vetoed alternatives into the Top-$K$ tier.

These results are consistent with Proposition~\ref{prop:rank_reversal}. Under linear aggregation, large predictive scores eventually compensate for a zero context score when prediction dominates the aggregation. In contrast, the conjunctive minimum and weighted geometric operators preserve the contextual exclusion independently of the magnitude of the predictive score, making the resulting prioritization substantially more robust to localized predictive overconfidence.

\subsection{Rank Stability Under imprecision in Context evaluation}
\label{subsect:noise}

This experiment evaluates how imprecision in the context score affects the stability of the final prioritization. Context scores derived from fuzzy rules or expert assessments are inherently imprecise, and small variations may alter the ordering of alternatives. 
Imprecision here is simulated perturbing the context scores using Gaussian with a standard deviation $\sigma$. Then, the rank stability is measured by Kendall's rank correlation, $\bar{\tau}(\sigma)$, computed between the baseline ranking ($\sigma=0$) and the ranking obtained after perturbation.

Figure~\ref{fig:noise_propagation} compares the linear ($A_L$) and minimum ($A_C^{\min}$) operators at $\lambda=0.35$. At $\sigma=0$, both operators reproduce the baseline ranking exactly ($\tau=1$ in every Monte Carlo trial). As the imprecision level increases, the ranking produced by the linear operator deteriorates more rapidly than that of the minimum operator. At $\sigma=0.5$, the average rank correlation decreases to approximately $0.40$ for $A_L$ (95\% CI $[0.38,0.43]$), compared with approximately $0.47$ for $A_C^{\min}$ (95\% CI $[0.44,0.50]$).

\begin{figure}[t]
\centering
\includegraphics[width=\textwidth]{figures/noise_propagation.pdf}
\caption{Rank stability under context uncertainty ($\lambda=0.35$, $N=30$ Monte Carlo trials). Solid curves represent Monte Carlo means and shaded regions denote 95\% confidence intervals. The vertical axis reports Kendall's rank correlation, $\bar{\tau}(\sigma)$, between the baseline ranking and the ranking obtained after perturbing the context scores with Gaussian noise of standard deviation $\sigma$.}
\label{fig:noise_propagation}
\end{figure}

This behavior is consistent with Proposition~\ref{prop:cont_homog_rank_preserv}. Under linear aggregation, every perturbation in the context score contributes directly to the priority score, so imprecision propagates throughout the ranking. In contrast, the minimum operator is affected only when the perturbed context score becomes smaller than the predictive score. Whenever $R_i \leq Q_i$, the aggregated score remains equal to $R_i$, making the prioritization insensitive to variations in the context score above that threshold.

Together with the previous experiments, these results show that the choice of aggregation operator influences not only the balance between prediction and context, but also the robustness of the resulting prioritization. While the linear operator performs well within a limited range of $\lambda$, the non-linear operators preserve contextual constraints and exhibit greater robustness to imprecision in the context score.

{\color{blue}
\section{Analysis in a Realistic Case Study}
\label{sect:real_example}

The previous experimental analysis isolated the behavior of the aggregation operators under controlled conditions. We now examine whether similar patterns are observed in a realistic CADEMAS-ML application using the employee attrition case study introduced by \citet{Godz2026}. The analysis considers the complete cohort of $n=100$ employees and a resource-constrained intervention policy selecting the Top-$K$ alternatives with $K=10$. Each employee is represented directly by a predictive score $R_i$ and a context score $Q_i$, obtained from the original CADEMAS-ML workflow.

The predictive score aggregates the outputs of the cooperative ADM units, whereas the context score represents the degree of alignment of an employee with the \textit{Digital Transformation} context. In this cohort, $R_i\in[0.001,0.989]$, $Q_i\in[0,0.73]$, and 35 employees receive $Q_i=0$, corresponding to contextually excluded cases. The observed attrition label in the original dataset is retained only for validation and does not participate in the prioritization process.

Five integration settings are compared: the linear operator $A_L$ with $\lambda=\{0.1,0.5,0.9\}$, together with the minimum ($A_C^{\min}$) and weighted geometric ($A_G$) operators. Particularly, for $A_G$ we employed $\lambda=0.5$. Table~\ref{tab:real_example} summarizes the Top-$K$ employees obtained under the five integration settings, whereas Figures~\ref{fig:real_example_plane} and~\ref{fig:real_example_bump} provide complementary visualizations of the employee distribution in the $(Q_i,R_i)$ space and the corresponding rank evolution across operators. Names are fictitious and are used only to facilitate the discussion of representative cases.

\begin{table}[H]
    \caption{Top-$K$ employees across the five integration settings for the 100-employee attrition cohort (Digital Transformation scenario, $K=10$). The first block lists the baseline Top-$K$ tier under $A_L$ with $\lambda=0.5$; the second block lists employees who enter the Top-$K$ tier under at least one alternative operator. In the baseline block, plain $P_i$ values denote Top-$K$ membership and bold values mark exclusion; in the second block, bold italic values mark entry under the corresponding operator.}
    \label{tab:real_example}
    \renewcommand{\arraystretch}{1.3}\scriptsize
    \begin{adjustwidth}{-\extralength}{0cm}
    \begin{tabularx}{\fulllength}{>{\raggedright\arraybackslash}p{2.4cm}|>{\raggedright\arraybackslash}l|R|R|R|R|R|R|R|C}
    \toprule
    \textbf{Group} & \textbf{Employee} & $R_i$ & $Q_i$ & $P_i^{(L,0.1)}$ & $P_i^{(L,0.5)}$ & $P_i^{(L,0.9)}$ & $P_i^{(C^{\min})}$ & $P_i^{(G)}$ & \textbf{Attrition} \\
    \midrule
    \multirow[m]{10}{*}{\shortstack{Baseline Top-$K$\\($A_L$, $\lambda=0.5$)}} & Laura Baker & 0.9359 & 0.6000 & 0.6336 & 0.7679 & 0.9023 & 0.6000 & 0.7493 & Yes \\
    & Andrew Wood & 0.9288 & 0.5667 & 0.6029 & 0.7477 & 0.8925 & 0.5667 & 0.7255 & Yes \\
    & Alice Brown & 0.9295 & 0.5600 & 0.5969 & 0.7447 & 0.8925 & 0.5600 & 0.7215 & Yes \\
    & Andrew Ross & 0.9302 & 0.4667 & 0.5130 & 0.6984 & \textbf{0.8839} & 0.4667 & 0.6589 & Yes \\
    & Caroline Bennett & 0.9368 & 0.4571 & 0.5051 & 0.6970 & \textbf{0.8888} & 0.4571 & 0.6544 & Yes \\
    & Jessica Wilson & 0.9444 & 0.4000 & \textbf{0.4544} & 0.6722 & \textbf{0.8900} & 0.4000 & 0.6146 & Yes \\
    & Noah Murphy & 0.9743 & 0.3600 & \textbf{0.4214} & 0.6672 & 0.9129 & 0.3600 & 0.5923 & Yes \\
    & Eleanor Foster & 0.9574 & 0.3333 & \textbf{0.3957} & 0.6454 & 0.8950 & 0.3333 & 0.5649 & Yes \\
    & Christian Bennett & 0.9690 & 0.3200 & \textbf{0.3849} & 0.6445 & 0.9041 & \textbf{0.3200} & \textbf{0.5569} & Yes \\
    & Madison Foster & 0.9193 & 0.3667 & \textbf{0.4219} & 0.6430 & \textbf{0.8641} & 0.3667 & 0.5806 & Yes \\
    \midrule
    \multirow[m]{10}{*}{\shortstack{Top-$K$ entrants under\\alternative operators}} & Leah Cooper & 0.9611 & 0.3000 & 0.3661 & 0.6306 & \textbf{\textit{0.8950}} & 0.3000 & 0.5370 & Yes \\
    & Amelia Davis & 0.9879 & 0.2667 & 0.3388 & 0.6273 & \textbf{\textit{0.9158}} & 0.2667 & 0.5133 & Yes \\ 
    & Lucy Scott & 0.8747 & 0.3667 & 0.4175 & 0.6207 & 0.8239 & \textbf{\textit{0.3667}} & \textbf{\textit{0.5663}} & Yes \\
    & Dominic Collins & 0.9774 & 0.2000 & 0.2777 & 0.5887 & \textbf{\textit{0.8996}} & 0.2000 & 0.4421 & Yes \\
    & Camila Phillips & 0.9893 & 0.0000 & 0.0989 & 0.4947 & \textbf{\textit{0.8904}} & 0.0000 & 0.0000 & Yes \\
    & Julia Smith & 0.0048 & 0.7333 & \textbf{\textit{0.6605}} & 0.3691 & 0.0776 & 0.0048 & 0.0592 & No \\ 
    & Caroline Adams & 0.0566 & 0.6000 & \textbf{\textit{0.5457}} & 0.3283 & 0.1109 & 0.0566 & 0.1842 & Yes \\
    & Caroline Roberts & 0.0203 & 0.6000 & \textbf{\textit{0.5420}} & 0.3102 & 0.0783 & 0.0203 & 0.1104 & No \\
    & Hunter Adams & 0.0110 & 0.5714 & \textbf{\textit{0.5154}} & 0.2912 & 0.0671 & 0.0110 & 0.0794 & No \\
    & Claire Bailey & 0.0282 & 0.5333 & \textbf{\textit{0.4828}} & 0.2808 & 0.0787 & 0.0282 & 0.1227 & No \\
    \bottomrule 
    \end{tabularx}
    \end{adjustwidth} 
    %\noindent{\footnotesize \textbf{Bold italic} $P_i$ values: employee enters the Top-$K$ tier under the corresponding operator (second block only). \textbf{Bold} $P_i$ values: employee in the baseline block is excluded from that Top-$K$ tier.}
    \end{table}

\begin{figure}[H]
\centering
\includegraphics[width=0.7\textwidth]{figures/real_example_plane.pdf}
\caption{Employee distribution in the predictive vs. context space. Grey points represent the complete cohort ($n=100$). Colored circles indicate the baseline Top-$K$ tier obtained with the linear operator ($A_L$, $\lambda=0.5$), whereas star markers denote employees entering the Top-$K$ tier under at least one aggregation setting.}
\label{fig:real_example_plane} 
\end{figure}

\begin{figure}[H]
\begin{adjustwidth}{-\extralength}{0cm}
\hfill
\includegraphics[width=0.8\fulllength]{figures/real_example_bump.pdf}
\end{adjustwidth}
\caption{Rank evolution of the baseline Top-$K$ employees across the five integration settings. Star markers denote employees entering the Top-$K$ tier under an alternative operator or parameter setting. For these additional employees, the numeric prefix in the label indicates their rank in the baseline configuration ($A_L$, $\lambda=0.5$).}
\label{fig:real_example_bump}
\end{figure} 

The results in Table~\ref{tab:real_example} show that decreasing $\lambda$ shifts the linear operator toward a progressively more context-driven prioritization. Employees with strong contextual alignment but only modest predictive scores move upward in the ranking, even when their estimated attrition risk is low. The most extreme example is Julia Smith, whose context score is the highest in the cohort ($Q_i=0.733$) but whose predictive score is only $R_i=0.005$. Under the baseline configuration ($A_L$, $\lambda=0.5$), she ranks 48th, yet under $A_L$ with $\lambda=0.1$ she becomes the highest-priority employee. Similar behavior is observed for Caroline Adams, Caroline Roberts, Hunter Adams, and Claire Bailey, all of whom enter the Top-$K$ exclusively under the context-oriented linear configuration. Figure~\ref{fig:real_example_plane} shows that these employees occupy the upper-left region of the $(Q_i,R_i)$ plane, where strong contextual support compensates for weak predictive evidence.

Increasing $\lambda$ produces the opposite effect. As predictive evidence becomes dominant, employees with high predicted attrition risk but weak contextual alignment progressively move toward the intervention tier. Table~\ref{tab:real_example} shows that Camila Phillips, who has the highest predictive score in the cohort ($R_i=0.989$) but a null context score ($Q_i=0$), rises from rank 26 in the baseline configuration to the Top-$K$ under $A_L$ with $\lambda=0.9$. Leah Cooper, Amelia Davis, and Dominic Collins exhibit the same tendency, although with non-zero contextual scores. Their locations in Figure~\ref{fig:real_example_plane} correspond to the lower-right portion of the score space, illustrating how highly predictive but weakly aligned employees become increasingly favored as the predictive component receives greater weight.

The non-linear operators exhibit markedly different behavior. As shown in Table~\ref{tab:real_example}, the conjunctive minimum operator never allows employees with $Q_i=0$ to enter the prioritized set, regardless of their predictive score. The weighted geometric mean behaves less restrictively but still penalizes severe imbalance between prediction and context, preventing contextually excluded employees from reaching the intervention tier. Consequently, both operators preserve the contextual constraints identified in the theoretical analysis, whereas the linear operator eventually overrides them for sufficiently large values of $\lambda$.

The minimum operator also introduces a characteristic behavior absent from the compensatory schemes. Because the aggregated score is determined by the limiting component, different employees frequently receive identical priority values, producing ties in the final ranking. This effect is visible in Table~\ref{tab:real_example}, where several employees share exactly the same aggregated score under $A_C^{\min}$, and becomes more apparent in the rank trajectories of Figure~\ref{fig:real_example_bump}, where groups of employees collapse onto the same ranking level. Such ties are a direct consequence of the conjunctive semantics of the minimum operator rather than an artifact of the dataset.

Finally, Figure~\ref{fig:real_example_bump} provides a global view of how the Top-$K$ tier evolves across integration settings. The baseline employees remain relatively stable under moderate changes of $\lambda$, whereas larger shifts in the integration semantics produce selective replacements rather than a complete reordering of the prioritized set. Employees entering the Top-$K$ originate from the extremes of the $(Q_i,R_i)$ plane shown in Figure~\ref{fig:real_example_plane}: highly contextual but weakly predictive employees are favored by context-oriented linear integration, whereas highly predictive but weakly contextual employees emerge only under prediction-oriented linear integration. In contrast, both non-linear operators substantially reduce these extreme transitions by enforcing stronger interaction between predictive evidence and contextual assessment.

Overall, the realistic case study confirms the theoretical and simulation results. Linear aggregation behaves as a continuously adjustable compromise between prediction and context, while the minimum and weighted geometric operators encode fundamentally different integration policies that cannot be reproduced by merely tuning $\lambda$. Consequently, the choice of aggregation operator has a direct and practically significant impact on which employees ultimately receive intervention priority.
}

\section{Conclusions and Future Directions}\label{sect:conclusions}

This paper investigated how the aggregation operator shapes the integration of predictive and contextual information within CADEMAS, a general framework for cooperative automated decision-making \citep{NovoaHernandez2026}. Rather than treating aggregation as a purely mathematical operation, we showed that different operators encode different decision semantics and therefore lead to different prioritization behaviors.

The theoretical analysis established the conditions under which contextual information can modify or preserve the predictive ordering of alternatives. The experimental study demonstrated that these properties have practical consequences. Linear aggregation provides a flexible trade-off between predictive and context scores, but may allow contextually excluded alternatives to enter the prioritized set when the predictive component dominates the aggregation. In contrast, the minimum and weighted geometric operators preserve contextual exclusions while exhibiting greater robustness to localized predictive overconfidence and to imprecision in the context scores. The employee attrition case study confirmed that these behaviors remain visible in a realistic CADEMAS-ML application, illustrating how the choice of aggregation operator directly influences the final intervention priorities.

These results suggest that the aggregation operator should be regarded as a fundamental design choice in cooperative decision-making systems rather than as an implementation detail. Different organizations may attach different meanings to the context score: in some applications it represents guidance that can be balanced against predictive evidence, whereas in others it represents requirements that should not be compensated by high predictive scores. Selecting an aggregation operator therefore amounts to selecting the decision semantics through which predictive evidence and contextual knowledge interact.

This work has several limitations. The simulation study relies on synthetic populations designed to isolate individual behaviors of the aggregation operators rather than reproduce the full complexity of real applications. The case study uses precomputed predictive and context scores from an existing CADEMAS-ML instantiation instead of retraining the underlying predictive models. In addition, the analysis considers a representative subset of aggregation operators and focuses on ranking-based prioritization under resource constraints, without addressing issues such as fairness, intervention costs, or temporal adaptation of predictive models.

Future work will extend this analysis in several directions. An immediate step is to evaluate the proposed methodology on different application domains and with larger real-world cohorts. It will also be valuable to study broader families of aggregation operators, including parameterized fuzzy connectives and MCDM aggregation schemes, to better understand how different decision semantics influence prioritization. Another promising direction is to incorporate application-specific utility and cost models so that aggregation strategies can be evaluated not only by their ranking behavior but also by their practical consequences. Finally, and in the context of CADEMAS-ML, it is clear that  aggregation strategies should be considered as configurable components, allowing domain experts to explore their impact on policy compliance, prioritization, and decision robustness during system design.

%%%%%%%%%%%%%%%%%%%%%%%%%%%%%%%%%%%%%%%%%%
\authorcontributions{All authors contributed equally to this work. All authors have read and agreed to the published version of the manuscript.}

\funding{This research was supported by the project ``Study, Analysis and Evaluation of Cooperative Automated Decision-Making Systems (CADEMAS)'', PID2023-146575NB-I00, funded by MCIU/AEI/I0.13039/501100011033 and by FSE+. The authors acknowledge support of the previous project.}

\institutionalreview{Not applicable.}

\informedconsent{Not applicable.} 

\dataavailability{Simulation code, figure scripts, and cohort data for the real attrition example are available in the project repository at \url{https://github.com/pnovoahe/cademas-mathematics} (\texttt{experiments/}; 100-case CADEMAS-ML cohort in \texttt{experiments/data/attrition100/}). The underlying IBM HR Analytics dataset is publicly available on Kaggle.} 

\acknowledgments{Mariia Godz has a pre-doctoral research contract associated with the project ``Study, Analysis and Evaluation of Cooperative Automated Decision-Making Systems (CADEMAS)'', PID2023-146575NB-I00, funded by MCIU/AEI/I0.13039/501100011033 and by FSE+. The authors acknowledge support of the previous project.}

\conflictsofinterest{The authors declare no conflicts of interest.}

%%%%%%%%%%%%%%%%%%%%%%%%%%%%%%%%%%%%%%%%%%
%% Optional
%\appendixtitles{no} % Leave argument "no" if all appendix headings stay EMPTY (then no dot is printed after "Appendix A"). If the appendix sections contain a heading then change the argument to "yes".
%\appendixstart
%\appendix
%\section[\appendixname~\thesection]{}
%\subsection[\appendixname~\thesubsection]{}

%The appendix is an optional section that can contain details and data supplemental to the main text---for example, explanations of experimental details that would disrupt the flow of the main text but nonetheless remain crucial to understanding and reproducing the research shown; figures of replicates for experiments of which representative data are shown in the main text can be added here if brief, or as Supplementary Materials. Mathematical proofs of results not central to the paper can be added as an appendix.

%\begin{table}[H] 
%\caption{This is a table caption.\label{tab5}}
%\newcolumntype{C}{>{\centering\arraybackslash}X}
%\begin{tabularx}{\textwidth}{CCC}
%\toprule
%\textbf{Title 1}	& \textbf{Title 2}	& \textbf{Title 3}\\
%\midrule
%Entry 1		& Data			& Data\\
%Entry 2		& Data			& Data\\
%\bottomrule
%\end{tabularx}
%\end{table}

%\section[\appendixname~\thesection]{}

%All appendix sections must be cited in the main text. In the appendices, figures, tables, etc., should be labeled starting with “A”, e.g., Figure A1, Figure A2, etc.

%%%%%%%%%%%%%%%%%%%%%%%%%%%%%%%%%%%%%%%%%%
%\isPreprints{}{% This command is only used for ``preprints''.
\begin{adjustwidth}{-\extralength}{0cm}
%} % If the paper is ``preprints'', please uncomment this parenthesis.
%\printendnotes[custom] % Un-comment to print a list of endnotes

\reftitle{References}
\bibliography{biblio.bib}

% If authors have biography, please use the format below
%\section*{Short Biography of Authors}
%\bio
%{\raisebox{-0.35cm}{\includegraphics[width=3.5cm,height=5.3cm,clip,keepaspectratio]{Definitions/author1.pdf}}}
%{\textbf{Firstname Lastname} Biography of first author}
%
%\bio
%{\raisebox{-0.35cm}{\includegraphics[width=3.5cm,height=5.3cm,clip,keepaspectratio]{Definitions/author2.jpg}}}
%{\textbf{Firstname Lastname} Biography of second author}

% For the MDPI journals use author-date citation, please follow the formatting guidelines on http://www.mdpi.com/authors/references
% To cite two works by the same author: \citeauthor{ref-journal-1a} (\citeyear{ref-journal-1a}, \citeyear{ref-journal-1b}). This produces: Whittaker (1967, 1975)
% To cite two works by the same author with specific pages: \citeauthor{ref-journal-3a} (\citeyear{ref-journal-3a}, p. 328; \citeyear{ref-journal-3b}, p.475). This produces: Wong (1999, p. 328; 2000, p. 475)

%%%%%%%%%%%%%%%%%%%%%%%%%%%%%%%%%%%%%%%%%%
%% for journal Sci
%\reviewreports{\\
%Reviewer 1 comments and authors’ response\\
%Reviewer 2 comments and authors’ response\\
%Reviewer 3 comments and authors’ response
%}
%%%%%%%%%%%%%%%%%%%%%%%%%%%%%%%%%%%%%%%%%%
\PublishersNote{}
%\isPreprints{}{% This command is only used for ``preprints''.
\end{adjustwidth}
%} % If the paper is ``preprints'', please uncomment this parenthesis.
\end{document}